\section*{Chapitre \ref{ch:b3:quantitative-geomorphology} --- Géomorphologie quantitative}

\begin{solution}{exo:b3:quantitative-geomorphology:1}
$a/\sqrt{\pi\kappa t} = 0.30 - 0.04 = 0.26$, donc $\kappa t = (2/0.26)^2/\pi =
\qty{18.8}{m^2}$, et $t = 18.8/(5\times10^{-3}) = \qty{3.8}{ka}$, à un facteur
deux près, dans un sens comme dans l'autre, compte tenu de la plage de $\kappa$.
\end{solution}

\begin{solution}{exo:b3:quantitative-geomorphology:2}
$E = \kappa\,|\nabla^2 z| = 4\times10^{-3}\times0.012 =
\qty{4.8e-5}{m/an}$, \qty{0.048}{mm/an}. Relief $E L^2/(2\kappa) =
0.012\times80^2/2 = \qty{38.4}{m}$.
\end{solution}

\begin{solution}{exo:b3:quantitative-geomorphology:3}
$E = 2\times10^{-6}\times(25\times10^6)^{0.5}\times0.01 = 2\times10^{-6}
\times5000\times0.01 = \qty{1.0e-4}{m/an}$, \qty{0.1}{mm/an}. Comme $E$ est
en \unit{m/an}, $A^m$ en $\unit{m}^{2m}$ et $S$ n'a pas d'unité, $K$ s'exprime en
$\unit{m}^{1-2m}\,\unit{an^{-1}}$~: \unit{an^{-1}} pour $m = 0.5$.
\end{solution}

\begin{solution}{exo:b3:quantitative-geomorphology:4}
(a) Fixe~: un lit dur raidit la rivière là où il affleure, et le ressaut
reste attaché au lit. (b) Ancré à la faille~: tant que la faille joue,
la rivière est raidie là où elle la traverse. (c) Mobile~: le soulèvement de la
côte a abaissé le \omterm{def:b3:quantitative-geomorphology:knickpoint}{niveau de base} de toutes les rivières à la fois~; les \omterm{def:b3:quantitative-geomorphology:knickpoint}{knickpoints}
migrent vers l'amont et, comme le prédit le théorème de l'onde cinématique pour $n
= 1$, elles se situent sur une même courbe de niveau.
\end{solution}

\begin{solution}{exo:b3:quantitative-geomorphology:5}
$\theta = m/n = 0.5$. $S = (U/K)^{1/2}A^{-1/2} = (10^{-4}/10^{-9})^{1/2}
\times(10^7)^{-1/2} = 316\times3.16\times10^{-4} = 0.10$. Doubler $U$
multiplie $k_s = (U/K)^{1/n}$ par $\sqrt2 = 1.41$~; pour $n = 1$ par 2~: plus
$n$ est grand, moins la raideur réagit au soulèvement.
\end{solution}

\begin{solution}{exo:b3:quantitative-geomorphology:6}
Avec $x$ mesuré vers l'amont à partir de l'exutoire, $\mathrm{d}x = -\mathrm{d}d$, donc $\chi(d)
= \int_d^L (A_0/c\,d'^{h})^{\theta}\,\mathrm{d}d' = A_0^{\theta} c^{-\theta}
\int_d^L d'^{-h\theta}\,\mathrm{d}d' = A_0^{\theta}c^{-\theta}(L^p -
d^p)/p$. Ici $p = 0.19$, $A_0^{\theta}c^{-\theta} = 10^{2.7}/2^{0.45}$,
et $\chi = \qty{2935}{m}$.
\end{solution}

\begin{solution}{exo:b3:quantitative-geomorphology:7}
$c = 2\times10^{-6}\times(10^8)^{0.5} = \qty{0.02}{m/an}$, \qty{20}{m/ka}~;
\qty{5}{km} à cette vitesse, il faut \qty{250}{ka}. L'aire, et avec elle la
célérité, diminue vers l'amont, donc le \omterm{def:b3:quantitative-geomorphology:knickpoint}{knickpoint} ralentit en remontant.
\end{solution}

\begin{solution}{exo:b3:quantitative-geomorphology:8}
$\bar e = 0.40\times600 = \qty{240}{m}$~; les sommets s'élèvent de $0.80\times240 =
\qty{192}{m}$~; la hauteur moyenne baisse de $240 - 192 = \qty{48}{m}$.
\end{solution}

\begin{solution}{exo:b3:quantitative-geomorphology:9}
$h = h_0\ln(P_0/E)$~: $0.5\ln 6 = \qty{0.90}{m}$ et $0.5\ln1.5 =
\qty{0.20}{m}$. À \qty{0.4}{mm/an}, l'érosion dépasse la production la plus rapide,
$P_0$~: le sol est décapé et la pente devient de la roche nue, son abaissement
est alors limité par l'altération.
\end{solution}

\begin{solution}{exo:b3:quantitative-geomorphology:10}
{\sloppy À l'\omterm{def:b3:quantitative-geomorphology:steady-profile}{état stationnaire} $\mathrm{d}q/\mathrm{d}x = E$ et $q = 0$ à la ligne de partage, donc
$q = E x$. Alors $\kappa S = E x\,(1 - S^2/S_c^2)$, une équation du second degré dont la
racine positive est $S = [-\kappa + \sqrt{\kappa^2 + 4E^2x^2/S_c^2}\,]/(2Ex/S_c^2)$. Pour
20, 50, 100 et \qty{200}{m}~: 0.36, 0.68, 0.89 et 1.03, contre 0.4, 1.0,
2.0 et 4.0 pour une diffusion linéaire ($S = Ex/\kappa$). La colline présente un
sommet arrondi, puis des flancs dont la pente tend vers $S_c$~:
des versants presque plans, avec un relief bien moindre que la parabole.\par}
\end{solution}

\begin{solution}{exo:b3:quantitative-geomorphology:11}
Par le théorème de l'onde cinématique, la perturbation atteint un point de coordonnée
$\chi$ après $\chi/(KA_0^m)$~; le réseau a entièrement réagi lorsqu'il a
atteint la dernière tête de chenal, de plus grand $\chi$. $K A_0^m = 10^{-6}\times
1000 = \qty{e-3}{m/an}$~: le cours principal répond en \qty{5.46}{Ma}, le réseau
en \qty{7.24}{Ma}, de l'ordre du million d'années attendu.
\end{solution}

\begin{solution}{exo:b3:quantitative-geomorphology:12}
Au-dessus de la confluence, la pente stationnaire de l'affluent est $S = (U/2K)A^{-\theta}$,
donc $\mathrm{d}z/\mathrm{d}\chi = (U/2K)A_0^{-\theta}$, la moitié de celle du cours principal.
Le \omterm{def:b3:quantitative-geomorphology:knickpoint}{knickpoint} atteint la confluence lorsque celui du cours principal est en $\chi_j$~; après
un temps supplémentaire $\Delta t$, celui du cours principal est passé en $\chi_k$ avec $\chi_k -
\chi_j = K A_0^m\,\Delta t$, tandis que celui de l'affluent, avec une célérité $2KA^m$,
est passé en $\chi$ avec $\chi - \chi_j = 2KA_0^m\,\Delta t = 2(\chi_k -
\chi_j)$.
\end{solution}

\begin{solution}{pb:b3:quantitative-geomorphology:1}
\textbf{1.} $4.5\times12000^{1.7} = \qty{3.87e7}{m^2}$, \qty{38.7}{km^2}~;
$4.5\times300^{1.7} = \qty{7.3e4}{m^2}$, \qty{0.073}{km^2}.

\textbf{2.} $\theta \approx -\log(0.007855/0.06641)/\log(37.08/0.388) =
0.928/1.980 = 0.47$.

\textbf{3.} Oui~: avec une dispersion d'environ \qty{10}{\percent} sur douze points,
$0.466$ ne diffère pas significativement de 0.5.

\textbf{4.} $0.06641\times\sqrt{388\,000} = \qty{41.4}{m}$~;
$0.007855\times\sqrt{3.708\times10^7} = \qty{47.8}{m}$. La moyenne géométrique,
\qty{45.1}{m}, estime $k_s$ pour $\theta = 0.5$.

\textbf{5.} Une pente est la différence de deux hauteurs, chacune incertaine
de quelques décimètres à quelques mètres, divisée par une courte distance~; sur les
tronçons à faible pente, l'erreur représente une grande part de la différence. $\chi$ intègre au lieu de
dériver.

\textbf{6.} $p = 1 - 1.7\times0.5 = 0.15$~; $\chi = (1000/\sqrt{4.5})\,
(12000^{0.15} - 9000^{0.15})/0.15 = \qty{543.1}{m}$.

\textbf{7.} $(193.4 - 170)/543.1 = 0.0431$~; $k_{sn} = 0.0431\times A_0^{0.5}
= \qty{43.1}{m}$.

\textbf{8.} L'estimation en chi~: elle repose sur des hauteurs mesurées sur \qty{3}{km} de
cours d'eau, et non sur douze pentes bruitées.

\textbf{9.} $K = U/k_{sn} = 4.3\times10^{-5}/43.1 = \qty{1.00e-6}{an^{-1}}$.

\textbf{10.} $10^{-6}\times\sqrt{3.87\times10^7}\times0.00691 =
\qty{4.30e-5}{m/an}$, \qty{43}{m/Ma}~: l'exutoire suit le rythme du
soulèvement.

\textbf{11.} À l'exutoire $10^{-6}\times6222 = \qty{6.2}{m/ka}$~; à la
cascade, $d = \qty{7.40}{km}$, $A = \qty{1.70e7}{m^2}$ et
\qty{4.1}{m/ka}.

\textbf{12.} $170 + 43\times(4.5)^{-0.5}(12000^{0.15} - 300^{0.15})/0.15 =
\qty{405}{m}$~: un dénivelé de \qty{235}{m}.

\textbf{13.} Pour $n = 1$, l'équation est une onde cinématique~: les points de part et
d'autre de la marche s'élèvent de $U$ par unité de temps, si bien que la marche conserve sa hauteur,
soit les \qty{25}{m} de la chute.

\textbf{14.} $\chi = (1000/\sqrt{4.5})(12000^{0.15} - 7400^{0.15})/0.15 =
\qty{899.4}{m}$.

\textbf{15.} $t = \chi/(KA_0^{0.5}) = 899.4/10^{-3} = \qty{899}{ka}$.

\textbf{16.} Pour $\chi = 899.4 + 300 = \qty{1199}{m}$~:
$d^{0.15} = 12000^{0.15} - 1199\times0.15\times\sqrt{4.5}/1000$, $d =
\qty{6.25}{km}$, soit \qty{5.75}{km} en amont de l'exutoire.

\textbf{17.} Au confluent ($d = \qty{9.5}{km}$), $\chi = \qty{442.8}{m}$~;
la cascade se situe là où le $\chi$ de l'affluent atteint \qty{899.4}{m}~:
$s^{0.15} = 5000^{0.15} - (899.4 - 442.8)\times0.15\times\sqrt{1.5}/1000$,
$s = \qty{4.27}{km}$, soit \qty{0.73}{km} en amont du confluent.

\textbf{18.} Leurs confluents sont à $\chi = 1954$ et \qty{999}{m}, au-delà de
\qty{899}{m}~: le \omterm{def:b3:quantitative-geomorphology:knickpoint}{knickpoint} ne les a pas encore dépassés, si bien que le tronc
y est encore à son ancien niveau.

\textbf{19.} $E = 5\times10^{-3}\times0.0086 = \qty{4.3e-5}{m/an}$,
\qty{43}{m/Ma}~: les collines au-dessus de la cascade s'érodent au taux de soulèvement, en
\omterm{def:b3:quantitative-geomorphology:steady-profile}{état stationnaire}.

\textbf{20.} $150^2/(5\times10^{-3}) = \qty{4.5}{Ma}$, soit cinq fois l'âge de
la cascade.

\textbf{21.} Leur pied a été abaissé de \qty{25}{m} avec le chenal~: leurs
parties basses sont désormais trop raides, peut-être en glissement, tandis que leurs
crêtes continuent de s'éroder au même taux. Une vague de raidissement remonte chaque versant~; après quelques
millions d'années, elles retrouvent la même parabole, abaissée de \qty{25}{m}.

\textbf{22.} $t = \chi k_{sn}/(U A_0^{0.5})$ varie comme $1/U$~: de
$899/1.2 = 749$ à $899/0.8 = \qty{1124}{ka}$.

\textbf{23.} Avec $n = 2$, la célérité croît avec la pente~: la face raide
se propagerait plus vite que les tronçons voisins et resterait vive, mais
sa vitesse et l'âge que l'on en tire différeraient.

\textbf{24.} Le \omterm{def:b3:quantitative-geomorphology:knickpoint}{niveau de base} de la rivière Varden s'est abaissé
il y a \textbf{\qty[text-series-to-math=true]{900}{ka}} (\qty{899}{ka} par le
calcul).
\end{solution}
